\scanpage{31}\mspage{30}

Pour les contrées non ensemblistes, il y a lieu de supposer ceci (à vérifier sur les exemples non ensemblistes à venir)

\margin{(?)}
% slanted
\margin{remplace \uline{Cont} 4}
\uline{Cont} $\substack{1\\4}$ Pour que deux figures \inserted{$F,\unsure{C|G}$} soient \margin{\rule{0.35pt}{4.5\baselineskip}}compatibles, il f. et s. qu'elles le soient ensemblistement, et que pour deux \deleted{strates dist.} strates ouvertes $U$ de $F$, $V$ de $\unsure{C|G}$, $U\ne V$, \deleted{on ait} et pour \deleted{tout} $x\in U$, $y\in V$, $\deleted{\unsure{\{x\}\{\}\{|\{x\}\}}\unsure{-|=}\{x,y\}}$ $\{\{x\},\{y\}\}$ soit une figure (i.e $\{\{x\}\}$ et $\{\{y\}\}$ soient compatibles

% slanted
\margin{\uline{NB} Les contrées ``ensemblistes'' ou ``contrées de pts'' sont celles qui satisfont $\mathrm{Cont}'_{4}$, et telles que toute partie de $\mathcal L$ à 2 él. est une 0-sphère i.e. deux lieux $+$ distincts y sont disjoints.}

\uline{NB} Pour tout $x\in \mathcal L$, on note $\unsure{F|E}_x$ la figure $\{\{x\}\}$, consistant en la seule strate $\{x\}$.

% slanted
\margin{\uline{NB} Un lieu $x$ de $\mathcal L$ est dit \uline{ponctuel} si si $\forall y\deleted{\ne}\in L$, $y\ne x$, \unsure{$y$} est disjoint de $x$. Ainsi, les contrées de points sont celles où tout lieu est un pt. !}

\deleted{Dire que les fig} Si $x,y\in \mathcal L$ dire que $F_x$, $F_y$ sont compatibles, signifie que $\{\deleted{\unsure{\{x\},\{y\}}}F_x,F_y\}$ est \deleted{une} une figure. C'est toujours vrai si $x=y$. \deleted{Si $x\ne y$, \unsure{on dit alors}} \inserted{On dit} que \deleted{$\{$} $x,y$ sont ``\uline{disjoints}'' s'ils sont distincts et compatibles \deleted{\unsure{i.e ssi} \unsure{$\{x$}} On dit aussi \inserted{alors} que $\{x,y\}\in \mathcal P_2(\mathcal L)$ est une 0-\uline{sphère} de $C$.

Deux parties \deleted{\unsure{de}} $A,B$ de $\mathcal L$ sont dit \uline{disjointes} (au sens de la contrée $C$ -- ou $C$-disjointes) si $\forall x\in A$, $\forall y\in B$, $x,y$ sont disjoints ($\Longrightarrow A\cap B=\varnothing$).
